% =========================================================================
% Figura (apuntes manuscritos, pág. 7; teoría, pág. II-7): cambio de origen
% del momento dipolar. r'' = r' - OO'.
% =========================================================================
\begin{figure}[ht]
\centering
\begin{tikzpicture}[x=1cm, y=1cm, >=Stealth]
  \draw[magacero, line width=0.7pt, fill=rosaclaro!55]
    plot[smooth cycle, tension=0.8] coordinates {(-1.3,1.6) (0.2,3.1) (2.2,3.4) (3.4,2.4) (3.0,0.9) (1.4,0.3) (-0.3,0.4)};
  \node[etiq, text=magprof] at (3.55,3.15) {$V'$};
  \coordinate (O) at (0,0);
  \coordinate (Op) at (-0.75,1.45);
  \coordinate (Q) at (1.9,2.5);
  % ejes del origen O
  \draw[eje] (O) -- (0,1.0) node[above, etiq, font=\footnotesize] {$z$};
  \draw[eje] (O) -- (4.3,0) node[right, etiq] {$y$};
  \draw[eje] (O) -- (-0.9,-0.75);
  \fill[black] (O) circle (1.4pt) node[below right=1pt, etiq] {$O$};
  \fill[black] (Op) circle (1.4pt) node[left=2pt, etiq] {$O'$};
  \draw[magprof, line width=0.5pt, fill=white] ($(Q)+(-0.12,-0.12)$) rectangle ($(Q)+(0.12,0.12)$);
  \node[etiq, font=\footnotesize, anchor=south west] at ($(Q)+(0.12,0.08)$) {$\dVp$};
  \draw[vect] (O) -- ($(Q)+(-0.08,-0.12)$) node[pos=0.55, right=1pt, etiq] {$\rp$};
  \draw[vectrojo] (Op) -- ($(Q)+(-0.13,-0.02)$) node[pos=0.5, above=2pt, etiq, text=rojoteo] {$\vc{r}\,''$};
  \draw[vectverde] (O) -- (Op) node[pos=0.3, left=6pt, etiq, text=verdeteo] {$\OOp$};
  \node[cuadro, anchor=west, font=\footnotesize, align=left] at (4.9,1.8)
     {$\vc{r}\,''=\rp-\OOp$\\[3pt] $\vp_{O'}=\vp_{O}-Q\,\OOp$};
\end{tikzpicture}
\caption{Cambio de origen. Un mismo elemento de carga tiene posición $\rp$ respecto de $O$ y $\vc{r}\,''$ respecto de $O'$. El momento dipolar cambia en $-Q\,\OOp$, de modo que
no depende del origen si la distribución es neutra (\cref{prop:p-origen}).}
\label{fig:origen}
\end{figure}
